% Generated by report_gen.py (ai-ee v1) - do not hand-edit.
\documentclass[11pt,a4paper]{article}
\usepackage[margin=2.2cm]{geometry}
\usepackage{graphicx}
\usepackage{booktabs}
\usepackage{longtable}
\usepackage{pdfpages}
\makeatletter
\@ifundefined{KV@Gin@artifact}{\define@key{Gin}{artifact}[]{}}{}
\makeatother
\usepackage{xcolor}
\usepackage[hidelinks]{hyperref}
\setcounter{tocdepth}{1}
\begin{document}
\sloppy
\begin{center}
{\LARGE\bfseries buck-5v3a --- Design Document}\\[6pt]
{\large ai-ee v1 pipeline}\\[2pt]
generated 2026-08-08 23:21:23
\end{center}
\tableofcontents

\section{Board and Run Metadata}
{\small
\begin{longtable}{lp{11cm}}
\toprule
\textbf{Field} & \textbf{Value} \\
\midrule
\endhead
board & buck-5v3a \\
workspace & \texttt{boards/buck-5v3a} \\
phase & P2 \\
gate status & no gates recorded yet \\
state created & 2026-08-08T22:10:35 \\
state updated & 2026-08-08T23:16:50 \\
\bottomrule
\end{longtable}
}
\section{Overview}
\subsection*{Brief: buck-5v3a}
Build a synchronous buck converter board.

\begin{itemize}
\item Input 7-18 V, nominal 12 V, via 2-pin 5.08 mm screw terminal.
\item Output 5 V at 3 A continuous, to a second 2-pin 5.08 mm screw terminal.
\item Include reverse-polarity protection on the input.
\item Ambient up to 50 C, no forced airflow.
\item Board size \textless{}= 50 x 40 mm.
\end{itemize}
(Verbatim user task: "Build a synchronous buck converter board. Input 7-18V, nominal 12V, via 2-pin 5.08mm screw terminal. Output 5V at 3A continuous to a second screw terminal. Include reverse-polarity protection on the input. Ambient up to 50C, no forced airflow. \textless{}=50x40mm.")

\subsection*{Stackup}
Chosen stackup: \textbf{`JLC04162H-7628A`} - JLCPCB impedance-controlled 4-layer, 1.6 mm nominal,
\section{Requirements}
\subsection*{Requirements: buck-5v3a}
\subsubsection*{1. Function}
Standalone synchronous step-down (buck) DC-DC converter board. It takes a DC input of 7-18 V (12 V nominal) on a 2-pin 5.08 mm screw terminal and produces a regulated 5 V output at 3 A continuous (15 W) on a second 2-pin 5.08 mm screw terminal, with reverse-polarity protection on the input so a swapped supply does not damage the board or the load. Power-only board: no MCU, no digital interface, no communication bus stated. Must run at up to 50 C ambient with no forced airflow, in a 50 x 40 mm outline.

\subsubsection*{2. Interfaces}
\begin{itemize}
\item Power input (J1): 2-pin screw terminal, 5.08 mm pitch (stated). DC only.
ASSUMED: through-hole terminal block, standard 5.08 mm 2-pole (e.g.
Phoenix MKDS / MSTB-style clone as stocked by LCSC); rated \textgreater{}= 10 A and
\textgreater{}= 250 V so it is not the limiting element at 2.5 A input. Polarity is
silkscreen-marked (+ / -), and the board tolerates the reverse case by
design (see Power).
\item Power output (J2): 2-pin screw terminal, 5.08 mm pitch (stated). 5 V,
3 A continuous. ASSUMED: same terminal family/part as J1 for BOM
consolidation, but physically separated and silk-labelled distinctly
(VIN vs 5V OUT) so the two cannot be confused in the field.
\item No other external interfaces stated: no USB, no RF, no header, no
comms bus, no external enable/on-off line, no fault output. Indicators
and test points are unspecified - see Open questions 6.
\end{itemize}
\subsubsection*{3. Power}
\begin{itemize}
\item Source: external DC supply, 7-18 V, 12 V nominal (stated). Source TYPE
is NOT stated (bench supply / AC-DC brick / battery / vehicle 12 V
system) and is safety-relevant - see Open questions 1. The pipeline
will not proceed on a guessed answer.
\item Input current (GUESS, for sizing only): 15 W out at an assumed 90-93\%
efficiency -\textgreater{} \textasciitilde{}16.1-16.7 W in. At 12 V nominal \textasciitilde{}1.4 A; at the 7 V
low-line corner \textasciitilde{}2.3-2.4 A. Input path (terminal, reverse-polarity
device, fuse if fitted, inductor-side copper) sized for \textgreater{}= 2.5 A
continuous.
\item Output rail: single 5 V rail, 3 A continuous = 15 W (stated). No other
rail stated - no auxiliary 3.3 V, no bias rail exposed. Accuracy and
ripple limits are not stated - see Open questions 3.
\item Topology: synchronous buck (stated) - integrated-FET regulator or
controller-plus-FETs is the architect's call at P1/P2, but synchronous
rectification is a requirement, not a preference (a non-synchronous
design is out of spec).
\item Reverse-polarity protection: REQUIRED on the input (stated). Method not
stated - see Open questions 4. GUESS for thermal budgeting: a series
P-channel MOSFET costs \textasciitilde{}50-100 mW at 2.4 A, a Schottky diode \textasciitilde{}1.0-1.2 W
at the same current, which is a meaningful share of the total loss
budget on a 50 C-ambient, no-airflow board.
\item Loss/thermal budget (GUESS): \textasciitilde{}1.1-1.7 W total board dissipation at
12 V in / 3 A out. With 50 C ambient, no airflow, and a 50 x 40 mm
outline, copper area and via stitching under the regulator are design
drivers, not afterthoughts. Peak component temperature target ASSUMED
as \textless{}= 105 C junction-equivalent hotspot (derated), unless the answer to
Open questions 9 imposes an enclosure that raises the internal ambient.
\item No battery, no charging circuit, and no energy storage function stated
on this board - PENDING confirmation that the SOURCE is not a battery
(Open questions 1).
\end{itemize}
\subsubsection*{4. Environment}
\begin{itemize}
\item Ambient: up to 50 C (stated).
\item Cooling: natural convection only, no forced airflow (stated). ASSUMED:
no heatsink and no thermal interface to a chassis unless the answer to
Open questions 9 says the board mounts to metal.
\item Minimum ambient not stated. ASSUMED: 0 C (indoor/benign). If the source
turns out to be automotive or outdoor, this becomes -40 C and changes
capacitor and inductor selection - covered by Open questions 1 and 9.
\item Enclosure, ingress (IP) rating, vibration/shock, humidity, conformal
coating: none stated - see Open questions 9.
\end{itemize}
\subsubsection*{5. Size \& mounting}
\begin{itemize}
\item Outline: 50 x 40 mm maximum (stated). Treated as a HARD cap - it binds
permanently at P5 board\_init and cannot be relaxed later without a
restart of the mechanical baseline. Smaller is allowed; larger is not.
\item Mounting holes: not stated - see Open questions 8.
\item Height limit: not stated - see Open questions 8. Relevant because a
15 W buck at 50 C ambient may want a taller shielded inductor and/or a
tall electrolytic input capacitor, and screw terminals themselves stand
\textasciitilde{}10-12 mm tall with wiring.
\item Connector edge placement: not stated. ASSUMED: both screw terminals sit
on board edges with their wire-entry openings facing outward so field
wiring does not cross the board; the architect picks which edges.
\end{itemize}
\subsubsection*{6. Quantity \& budget}
\begin{itemize}
\item Build quantity: not stated - see Open questions 10. ASSUMED default of
a small prototype run (qty 5) for costing purposes only.
\item Target unit cost: not stated - see Open questions 10. ASSUMED: no hard
per-unit cap; cost minimized by preferring JLC Basic parts and an
economy fab/assembly tier, rather than hitting a fixed number.
\end{itemize}
\subsubsection*{7. Assembly}
\begin{itemize}
\item Not stated - see Open questions 11. ASSUMED: JLCPCB PCBA, single-sided
(top) SMT assembly, which suits a power board where the bottom copper
is wanted as an unbroken thermal/return plane.
\item The two 5.08 mm screw terminals are through-hole parts. ASSUMED: fitted
by JLCPCB through-hole assembly if offered at acceptable cost for the
chosen part, otherwise hand-soldered after SMT - either is acceptable
and does not change the schematic or the footprint.
\item No stated restriction on package sizes or on Extended/Basic library
parts; ASSUMED preference for JLC Basic where a Basic part meets the
electrical requirement, with Extended allowed for the regulator IC and
the power inductor (Basic stock rarely covers 3 A-class buck parts).
\end{itemize}
\subsubsection*{8. Compliance/safety flags}
Two flags apply; one is pending an answer.

\begin{itemize}
\item High current (\textgreater{}3 A): APPLIES at the boundary - 3 A continuous output is
the stated rating, and any specified peak/inrush allowance pushes the
board above 3 A (Open questions 2). Consequences carried forward:
trace/copper sizing with a defined temperature rise, thermal relief and
via stitching, input/output capacitor RMS ripple current ratings,
connector and fuse ratings, and a defined short-circuit behavior
(Open questions 5).
\item \textgreater{}30 V: DOES NOT APPLY at the stated steady-state maximum of 18 V, but
the input transient environment is unknown. If the source is a vehicle
12 V system (Open questions 1), load-dump and ISO 7637-2 transients can
exceed 40 V and this flag flips to APPLIES, forcing input clamping
(TVS), a higher-voltage regulator, and higher-voltage input capacitors.
Not guessed - asked.
\item Batteries: UNKNOWN, PENDING Open questions 1. The brief gives a voltage
window (7-18 V) that is equally consistent with a bench supply, an
AC-DC adapter, a 12 V lead-acid battery, or a 4S Li-ion pack. No
charging function is requested on this board, but battery-source status
changes reverse-polarity expectations (a mis-clipped battery is a
sustained, low-impedance reverse fault, not a brief supply glitch),
deep-discharge/UVLO behavior, and the minimum ambient. Per the P0
rules, this is asked, never assumed.
\item Mains voltage: does not apply - the board sees DC only, and generating
the DC input is out of scope for this board.
\item Motors: does not apply - no motor drive stated. If the 5 V load is
actually a motor or an inductive/stalling load, the peak-current answer
(Open questions 2) must reflect it.
\item RF transmit: does not apply - no radio function.
\end{itemize}
\subsubsection*{9. Open questions}
STATUS: ALL RESOLVED 2026-08-08 by the user (Q1, Q2, Q4, Q11 answered explicitly; Q3, Q5-Q10 accepted at their stated DEFAULT). The binding answers are in section 10 - where section 10 and this section differ, section 10 wins.

1. What exactly powers the input? Pick one: (a) bench power supply, (b) AC-DC wall adapter or brick, (c) a battery - and if so, which chemistry (lead-acid, LiFePO4, Li-ion pack, other), (d) a vehicle / automotive 12 V system, or (e) something else. This is the one safety-relevant question with NO default: it decides whether the board needs load-dump/surge clamping (up to 40 V+), a wider temperature range, and battery-specific reverse-connection and under-voltage behavior. Please answer explicitly. 2. Is 3 A the absolute maximum output current, or must the board also survive brief peaks above it (for example a motor or capacitive load drawing 4-5 A for tens of milliseconds at startup)? If there are peaks, how much current and for how long? DEFAULT: 3 A continuous is the maximum; no specified peak above it, and the regulator's own current limit provides the headroom. 3. How tight does the 5 V output need to be? Two numbers: DC accuracy (how far from exactly 5.00 V is acceptable) and ripple (how much high-frequency wobble riding on the 5 V is acceptable). DEFAULT: +/-3\% DC accuracy (4.85-5.15 V) over the full line and load range, and \textless{}= 50 mV peak-to-peak output ripple. 4. Which style of reverse-polarity protection do you want? A series P-channel MOSFET wastes very little power (tens of mW) but adds a few parts; a Schottky diode is one cheap part but burns roughly 1 W as heat, which is significant on a fanless board at 50 C ambient. DEFAULT: P-channel MOSFET (low loss - the right call for this thermal budget). 5. Should the board include its own input fuse, and what should happen on an output short-circuit? Options for the fuse: a one-shot cartridge or SMD fuse (must be replaced after a fault), a resettable polyfuse (recovers by itself), or none (rely on the upstream supply's limit). DEFAULT: fit one non-resettable SMD fuse rated \textasciitilde{}4 A on the input, and rely on the regulator IC's built-in hiccup-mode current limit and thermal shutdown for output short-circuit and overload protection. 6. What local indicators and test access do you want? Options: a power LED on the 5 V output, test points for probing, an external enable/on-off input, a power-good signal brought out. DEFAULT: one green power LED on the 5 V output plus three test points (VIN, 5V, GND); no external enable input and no power-good output (nothing in the brief needs to talk to this board). 7. Does this board have to pass a formal EMC/emissions test (for example CISPR 32 for an information-technology product, or CISPR 25 if it goes in a vehicle), or is it an internal/prototype board where good layout practice is enough? DEFAULT: no formal EMC test campaign; design with power-loop-tight layout and normal input filtering, and leave room for an optional input filter if you later need it. 8. Mechanically: does the board need mounting holes, and is there a height limit above the PCB? If it bolts into something, what hole size and pattern? DEFAULT: four M3 holes (3.2 mm drill) inset from the corners of the 50 x 40 mm outline, and a 15 mm maximum component height (the screw terminals themselves are about 10-12 mm tall). 9. Where does the board live? Specifically: open on a bench, inside a plastic or metal enclosure, exposed to dust/water (any IP rating), subject to vibration, or needing conformal coating? Also, what is the coldest ambient it will see? DEFAULT: open board, indoors, no enclosure, no ingress rating, no vibration requirement, no conformal coating, and a 0 C minimum ambient. 10. How many boards will be built, and is there a hard per-unit cost target? DEFAULT: a prototype run of 5 boards, with no hard cost cap - cost is minimized by preferring JLC Basic parts and an economy fab tier. 11. How should it be assembled? Options: JLCPCB machine assembly of the surface-mount parts only (you hand-solder the two screw terminals), JLCPCB assembly including the through-hole screw terminals, or full hand assembly. Also: is it acceptable for all surface-mount parts to sit on the top side only? DEFAULT: JLCPCB PCBA, top-side surface-mount assembly only, with the two screw terminals added by JLCPCB through-hole assembly if the cost is reasonable and hand-soldered otherwise; bottom side kept clear as a thermal/ground plane.

\subsubsection*{10. Answers (binding, user-confirmed 2026-08-08)}
A1. SOURCE: bench power supply or AC-DC adapter/brick. NOT a battery, NOT automotive. Consequences: the \textgreater{}30 V flag stays CLOSED - no load-dump clamping campaign, no ISO 7637-2, no -40 C rating. Input steady-state max stays 18 V; parts rated 25-40 V with a modest input TVS for hot- plug/inductive-cable ringing are sufficient. Battery flag CLOSED: no UVLO/deep-discharge behavior required, and the reverse-connect fault is supply-current-limited rather than a pack short. A2. PEAK LOAD: 3 A is the absolute maximum. No peak allowance above 3 A. Size the inductor saturation and current limit for the regulator's own limit (target \textasciitilde{}4-5 A), not for an external burst. A3. OUTPUT TIGHTNESS (default accepted): +/-3 \% DC (4.85-5.15 V) over the full 7-18 V line and 0-3 A load range; \textless{}= 50 mV pk-pk output ripple. A4. REVERSE-POLARITY: series P-channel MOSFET in the input path (gate clamped, e.g. gate resistor + zener). Schottky rejected on thermal grounds. Target conduction loss \textless{}= 100 mW at the 2.4 A low-line input current. A5. FUSE / SHORT-CIRCUIT (default accepted): one non-resettable SMD fuse on the input, \textasciitilde{}4 A rating; output short-circuit and overload handled by the regulator's hiccup current limit plus thermal shutdown. A6. INDICATORS / TEST ACCESS (default accepted): one green power LED on the 5 V output; three test points (VIN, 5V, GND). No external enable input, no power-good output. A7. EMC (default accepted): no formal test campaign. Tight power-loop layout and normal input filtering; leave room for an optional input filter. A8. MECHANICAL (default accepted): four M3 (3.2 mm) holes inset from the corners; 15 mm maximum component height above the PCB. A9. ENVIRONMENT (default accepted): open board, indoors, no enclosure, no IP rating, no vibration requirement, no conformal coating, 0 C minimum ambient. 50 C maximum ambient and natural convection stand. A10. QUANTITY / COST (default accepted): prototype run of 5. No hard per-unit cap; prefer JLC Basic parts and an economy fab tier. A11. ASSEMBLY: JLCPCB PCBA, TOP-SIDE SMT ONLY. The two 5.08 mm screw terminals go to JLC through-hole assembly if reasonably priced, and are hand-soldered otherwise - either way the schematic and footprint are unchanged. Bottom side stays clear of SMT parts and is used as a thermal/return plane. Prefer Basic/Preferred parts; Extended allowed for the regulator IC and the power inductor.

\section{Architecture}
\subsection*{Decisions of record (state.json)}
{\small
\begin{longtable}{lp{5.2cm}p{7.2cm}}
\toprule
\textbf{Phase} & \textbf{What} & \textbf{Why} \\
\midrule
\endhead
P0 & Input source is bench supply / AC-DC adapter; NOT battery, NOT automotive & User answer A1: closes the \textgreater{}30V and battery safety flags; 18V steady-state max stands, no load-dump campaign, 0C min ambient \\
P0 & 3 A is the absolute output max; no peak allowance above 3 A & User answer A2: inductor saturation and current limit sized to regulator limit (\textasciitilde{}4-5 A), not to an external burst \\
P0 & Reverse-polarity protection = series P-channel MOSFET, target \textless{}=100 mW at 2.4 A & User answer A4: Schottky rejected on thermal grounds (50C ambient, no airflow, 50x40mm) \\
P0 & JLCPCB PCBA top-side SMT only; screw terminals THT (JLC if priced sanely, else hand); bottom side reserved as thermal/return plane & User answer A11 \\
P0 & Defaults accepted for output tolerance (+/-3\%, \textless{}=50mV ripple), 4A input fuse + hiccup limit, power LED + 3 test points, no formal EMC, 4x M3 holes + 15mm height, indoor open board 0-50C, qty 5 proto, JLC Basic preference & User accepted requirements-analyst defaults Q3,Q5-Q10 in the P0 batch \\
P2 & Stackup JLC04162H-7628A: 4 layers, 2 oz outer / 0.5 oz inner, 1.6 mm; outline 50x40 mm (the full cap) & check\_thermal on the real part: AP63356Q 0.88 W at the 7 V corner -\textgreater{} Tj 95 C on 4L (pass) vs 115 C on 2L (fail). 45x35 does not fit \textasciitilde{}557 mm2 of parts plus the GND pour and is the same JLC price tier \\
P2 & Lead regulator AP63356QZV-7 (PWM-only sibling), not AP63357Q & PFM light-load ripple threatens the \textless{}=50 mV pk-pk spec which has no stated minimum load; light-load efficiency is not a requirement here \\
P2 & CONFLICT RULED: datasheet theta\_JA 25 C/W rejected in favour of the repo calibrated 51 C/W (4L) & Datasheet figure is a JEDEC 8710 mm2 coupon; this board gives \textasciitilde{}2000 mm2, and 51 C/W is what the P8 verify gate applies - using the datasheet number would pass P2 and fail P8 \\
P2 & power.json's '\textless{}=90 mohm RDS(on)' part filter DELETED; loss table re-derived at 1.48 W / 91.0 \% / Iin 2.35 A at 7 V & The filter was derived from generic part classes and would have rejected the entire real shortlist including the part that measurably passes \\
P2 & ONE FLAT SCHEMATIC SHEET, no hierarchy & One function, one rail, 25 parts; also dodges the /\textless{}sheet\textgreater{}/\textless{}LABEL\textgreater{} net-prefix trap. Deviation from the default hierarchical plan, recorded per full-run.md \\
\bottomrule
\end{longtable}
}
\subsection*{Sheet plan}
\subsection*{buck-5v3a - schematic sheet plan}
\subsubsection*{1. ONE FLAT SHEET - and why that is the right answer, not laziness}
{*}{*}No hierarchy. One root sheet, \texttt{buck-5v3a.kicad\_sch}, containing all 25 placed parts.{*}{*} This is a deliberate deviation from the pipeline's usual hierarchical default, recorded here so it is not mistaken for an omission:

\begin{itemize}
\item The board has \textbf{one function} (7-18 V -\textgreater{} 5 V/3 A) and \textbf{one rail}. There is no
second function to isolate, no repeated block to instance, no MCU/peripheral split.
\item Hierarchy costs a real thing here: KiCad names a net from a hierarchical label as
\textbf{`/\textless{}sheet\textgreater{}/\textless{}LABEL\textgreater{}`}, so every net in \texttt{constraints.json} would acquire a sheet
prefix, and {*}{*}a net name that does not match is silently no-op'd by every P5-P8
consumer{*}{*} - the trap that cost the lumina-par run a P4 amendment. A flat sheet
makes local labels come out as plain \texttt{/SW}, \texttt{/FB}, and removes that failure mode
entirely.
\item 25 parts fit on one A4 sheet with room to lay it out in signal-flow order
(left to right: J1 -\textgreater{} F1 -\textgreater{} Q1 -\textgreater{} Cin/U1 -\textgreater{} L1 -\textgreater{} Cout -\textgreater{} J2), which is also the
board floorplan. One sheet, one floorplan, one reading order.
\end{itemize}
The threshold for revisiting this: if P3/P4 adds an input filter stage or a second rail, split at that point - not before.

\begin{quote}\small\ttfamily\raggedright
\textbar{} Sheet \textbar{} Blocks (`blocks.md`) \textbar{} Refdes block \textbar{} `\#PWR` base \textbar{}\\
\textbar{}---\textbar{}---\textbar{}---\textbar{}---\textbar{}\\
\textbar{} {*}{*}root{*}{*} (flat) \textbar{} B1-B6, all of them \textbar{} {*}{*}1-99{*}{*} per prefix \textbar{} {*}{*}100{*}{*} (`\#PWR0100`+) \textbar{}
\end{quote}
Refdes uniqueness is structural on a single sheet. \texttt{pwr\_base = 100} is still declared explicitly so that a later hierarchy split (see threshold above) starts from a convention rather than colliding with \texttt{\#PWR01}.

\subsubsection*{2. Refdes assignment - BINDING ON P4}
\texttt{constraints.json} (\texttt{placement.groups}, \texttt{placement.separation}, \texttt{thermal.ref}) references these refdes \textbf{by name}. A rename silently drops the constraint that mentions it - \texttt{place\_anneal} collects unknown separation refs but a dropped grouping just produces a worse placement (S14 lesson). If P4 must rename, update \texttt{constraints.json} in the same commit.

\begin{quote}\small\ttfamily\raggedright
\textbar{} Ref \textbar{} Part / class \textbar{} Block \textbar{} Notes \textbar{}\\
\textbar{}---\textbar{}---\textbar{}---\textbar{}---\textbar{}\\
\textbar{} {*}{*}J1{*}{*} \textbar{} 5.08 mm 2P THT screw terminal, WJ500V-5.08-2P class \textbar{} B1 \textbar{} left edge, wire entry outward, silk `VIN 7-18V` \textbar{}\\
\textbar{} {*}{*}F1{*}{*} \textbar{} 4 A 1206 one-shot fuse, 1206T4A63V class \textbar{} B1 \textbar{} coolest corner, next to J1 \textbar{}\\
\textbar{} {*}{*}Q1{*}{*} \textbar{} P-FET, AO4407A class SO-8 (AOD403 DPAK alt) \textbar{} B2 \textbar{} source `/VIN\_FUSED`, drain `+VIN` \textbar{}\\
\textbar{} {*}{*}R6{*}{*} \textbar{} 10 kohm gate pull-down \textbar{} B2 \textbar{} NOT 100-330 ohm - `decisions.md` D6 \textbar{}\\
\textbar{} {*}{*}D3{*}{*} \textbar{} zener 12-15 V, Vgs clamp \textbar{} B2 \textbar{} Vz \textless{}= Vgs(max) - 5 V of the chosen Q1 \textbar{}\\
\textbar{} {*}{*}D2{*}{*} \textbar{} TVS, SMBJ20A class (20 V standoff / 32.4 V clamp) \textbar{} B3 \textbar{} on `+VIN`, AFTER Q1 \textbar{}\\
\textbar{} {*}{*}C1, C2{*}{*} \textbar{} 10 uF 50 V X7R 1210 \textbar{} B3 \textbar{} \textgreater{}= 1.5 A RMS at 450 kHz; X5R DISQUALIFIED \textbar{}\\
\textbar{} {*}{*}C3{*}{*} \textbar{} 100 nF 50 V 0603 \textbar{} B3 \textbar{} at the VIN pin, same layer, shortest loop \textbar{}\\
\textbar{} {*}{*}U1{*}{*} \textbar{} {*}{*}AP63356QZV-7{*}{*}, VDFN-13 3x2 \textbar{} B4 \textbar{} PWM-only sibling - `decisions.md` D3 \textbar{}\\
\textbar{} {*}{*}C4{*}{*} \textbar{} 100 nF bootstrap (BST-SW) \textbar{} B4 \textbar{} required, not optional \textbar{}\\
\textbar{} {*}{*}R1, R2{*}{*} \textbar{} FB divider 158k / 30.1k, {*}{*}0.5 \% or better{*}{*} \textbar{} B4 \textbar{} Vout = 0.8 x (1 + R1/R2) = 5.00 V \textbar{}\\
\textbar{} {*}{*}R3, R4{*}{*} \textbar{} EN/UVLO divider 86.6k / 20k, 1 \% \textbar{} B4 \textbar{} \textasciitilde{}6.2 V rising / \textasciitilde{}5.3 V falling \textbar{}\\
\textbar{} {*}{*}L1{*}{*} \textbar{} 6.8 uH, Isat \textgreater{}= 6 A, DCR \textless{}= 30 mohm, shielded, 125 C \textbar{} B5 \textbar{} {*}{*}P3 must sweep 6.8 uH - the P1 shortlist has none{*}{*} \textbar{}\\
\textbar{} {*}{*}C5, C6{*}{*} \textbar{} 22 uF 25 V X7R 1210 \textbar{} B5 \textbar{} \textasciitilde{}30 uF effective after 5 V DC bias \textbar{}\\
\textbar{} {*}{*}C7{*}{*} \textbar{} 100 uF polymer \textbar{} B6 \textbar{} {*}{*}RESERVED refdes, NOT FITTED{*}{*} (H1 open 1) \textbar{}\\
\textbar{} {*}{*}J2{*}{*} \textbar{} same terminal family as J1 \textbar{} B6 \textbar{} right edge, silk `5V 3A` \textbar{}\\
\textbar{} {*}{*}D1, R5{*}{*} \textbar{} green LED + 1 kohm \textbar{} B6 \textbar{} 1.9 mA \textbar{}\\
\textbar{} {*}{*}TP1, TP2, TP3{*}{*} \textbar{} test points \textbar{} B6 \textbar{} `+VIN`, `+5V`, `GND` \textbar{}\\
\textbar{} {*}{*}H1-H4{*}{*} \textbar{} M3 mounting holes \textbar{} - \textbar{} created by `board\_init --mounting-holes 4`, board-only (not in the schematic, excluded from BOM/CPL/parity) \textbar{}
\end{quote}
Prefix blocks for anything P3/P4 adds: \textbf{C8+, R7+, D4+, TP4+}. Do not renumber existing refdes to close gaps.

\subsubsection*{3. Net-naming contract - BINDING ON P4}
These are the \textbf{canonical names}. \texttt{constraints.json} is written against them, and a mismatch is a silent no-op in \texttt{check\_current}, \texttt{check\_thermal}, \texttt{rules\_gen}, \texttt{planes\_gen} and \texttt{route\_critical}.

\begin{quote}\small\ttfamily\raggedright
\textbar{} Net \textbar{} Type \textbar{} From -\textgreater{} to \textbar{} Declared in `constraints.json` \textbar{}\\
\textbar{}---\textbar{}---\textbar{}---\textbar{}---\textbar{}\\
\textbar{} `/VIN\_RAW` \textbar{} local label \textbar{} J1.1 -\textgreater{} F1 \textbar{} `power` 3.0 A, `pdn: false` \textbar{}\\
\textbar{} `/VIN\_FUSED` \textbar{} local label \textbar{} F1 -\textgreater{} Q1 source \textbar{} `power` 3.0 A, `pdn: false` \textbar{}\\
\textbar{} `+VIN` \textbar{} {*}{*}power symbol{*}{*} \textbar{} Q1 drain -\textgreater{} D2, C1-C3, U1.VIN, R3 \textbar{} `power` 3.0 A \textbar{}\\
\textbar{} `/VIN\_GATE` \textbar{} local label \textbar{} Q1 gate -\textgreater{} R6 -\textgreater{} D3 \textbar{} not declared (uA) \textbar{}\\
\textbar{} `/EN` \textbar{} local label \textbar{} R3/R4 midpoint -\textgreater{} U1.EN \textbar{} not declared (uA) \textbar{}\\
\textbar{} `/BST` \textbar{} local label \textbar{} C4 -\textgreater{} U1.BST \textbar{} not declared \textbar{}\\
\textbar{} `/SW` \textbar{} local label \textbar{} U1.SW -\textgreater{} L1, C4 low side \textbar{} `power` 3.6 A dT 20, `pdn: false` \textbar{}\\
\textbar{} `+5V` \textbar{} {*}{*}power symbol{*}{*} \textbar{} L1 -\textgreater{} C5, C6, J2, R5, R1 top \textbar{} `power` 3.6 A \textbar{}\\
\textbar{} `/FB` \textbar{} local label \textbar{} R1/R2 midpoint -\textgreater{} U1.FB \textbar{} not declared \textbar{}\\
\textbar{} `/LED\_A` \textbar{} local label \textbar{} R5 -\textgreater{} D1 anode \textbar{} not declared \textbar{}\\
\textbar{} `GND` \textbar{} {*}{*}power symbol{*}{*} \textbar{} everything \textbar{} `power` 3.6 A, `plane\_fed`, `pdn: false` \textbar{}
\end{quote}
Rules P4 must honour:

\begin{itemize}
\item \textbf{Power nets are bare global power symbols} (\texttt{+VIN}, \texttt{+5V}, \texttt{GND}); everything
else is a \textbf{root-sheet local label}, which on a flat sheet yields \texttt{/NAME} with no
prefix. Do not use hierarchical labels - there is no hierarchy.
\item \textbf{A power symbol's net name comes from its VALUE field, not its library pin name}
(LEARNINGS 2026-07-28). A \texttt{+VIN} symbol cloned from \texttt{+12V} and not re-valued will
produce the net \texttt{+12V} and silently orphan every \texttt{+VIN} constraint.
\item \texttt{/FB} senses at the \textbf{Cout node}, not at the inductor, and is routed away from
\texttt{/SW} and L1 (DS41948 Figure 47, buck.md s4).
\item U1's \texttt{COMP} pin is \textbf{tied to GND} (internal compensation). \texttt{PG} is left
unconnected; mark it with a no-connect flag so ERC stays clean. Pin 7 is \texttt{NC} -
leave it floating, do not tie it to anything.
\item \texttt{TP1} is on \textbf{`+VIN`} (post-protection), not on \texttt{/VIN\_RAW}.
\end{itemize}
\subsubsection*{4. Placement groups the sheet plan implies}
These become \texttt{placement.groups} in \texttt{constraints.json} and drive P6 clustering. They are listed here because the sheet layout should mirror them - a schematic drawn in this order makes the placement obvious.

1. \textbf{`hotloop`} - anchor \textbf{U1}, members \textbf{C1, C2, C3, C4}. The Cin -\textgreater{} VIN -\textgreater{} GND loop is the shortest loop on the board and outranks every other placement preference (buck.md s3, DS41948 rule 2). C4 (BST) sits directly against the BST/SW pins. 2. \textbf{`fb`} - anchor \textbf{U1}, members \textbf{R1, R2} plus \textbf{R3, R4}. FB network as close to the FB pin as possible (DS41948 rule 5); the EN divider rides along because it is the same small-signal corner of the part. 3. \textbf{`output`} - anchor \textbf{L1}, members \textbf{C5, C6}. Inductor as close to SW as possible (rule 3), output caps' GND return short (rule 4). 4. \textbf{`rpp`} - anchor \textbf{Q1}, members \textbf{R6, D3, D2}. The gate network must be at the FET; the TVS lands on the same \texttt{+VIN} node.

Separations (centroid-to-centroid, \texttt{place\_anneal} cost terms): \textbf{U1 \textless{}-\textgreater{} L1 \textgreater{}= 8 mm} (\textasciitilde{}2.5 mm body gap for 3x2 and 8x8 bodies) and \textbf{U1 \textless{}-\textgreater{} F1 \textgreater{}= 15 mm} (F1 in the cool corner by J1).

\subsubsection*{5. Area P6 must leave clear (no footprint, no constraint - a placement instruction)}
\begin{itemize}
\item \textbf{\textasciitilde{}40 mm\textasciicircum{}2 beside J2 for C7}, the reserved 100 uF polymer (D6.3 x 5.8 mm class),
so that a "yes" to H1 open question 1 is a BOM change and not a re-layout.
\item \textbf{\textasciitilde{}60 mm\textasciicircum{}2 between J1 and F1} for the optional input filter A7 asks room for. No
unpopulated footprints are added: a DNF part on a JLC PCBA BOM buys confusion, and
the requirement is for \emph{room}, not for parts.
\end{itemize}
Full architecture narrative (not inlined; see the artifact index): \texttt{architecture/blocks.md}, \texttt{architecture/decisions.md}, \texttt{architecture/power\_tree.md}, \texttt{architecture/stackup.md}.
\section{Schematic}
\emph{Pending --- produced at P4.}
\section{Layout}
\emph{Pending --- produced at P6.}
\section{Verification}
\emph{Pending --- produced at P8.}
\section{DFM and Fabrication}
\emph{Pending --- produced at P9.}
\section{Run Record (Appendix)}
\subsection*{Phase digests}
\paragraph*{P0}
\begin{quote}\small\ttfamily\raggedright
\# P0 Intake - digest\\
~\\
- buck-5v3a: 7-18 V (12 nom) -\textgreater{} 5 V / 3 A cont (15 W), screw terminals both\\
  ends, RPP on input, 50 C ambient natural convection, 50x40 mm HARD cap.\\
- requirements.md written; check\_requirements.py pass (0 violations, 1-9).\\
- 11 open questions batched to user, ALL resolved; binding answers = section 10.\\
- Safety flags CLOSED by A1 (bench/adapter source): no load dump, no battery\\
  case, 0-50 C. A2 caps output at 3 A, no peak allowance. High-current\\
  consequences (copper, cap RMS, thermal vias) carried forward.\\
- User calls: P-FET RPP (not Schottky, thermal); JLC PCBA top-side SMT only,\\
  bottom = thermal/return plane; 4 A input fuse; LED + 3 TPs; 4x M3, 15 mm\\
  height; qty 5 proto; JLC Basic preference.\\
- Deviation: requirements-analyst ran opus/high, not fable/high - Fable 5 out\\
  of usage credits.\\
- Next P1: power-architect + scout(regulator) + scout(magnetics/protection)\\
  in parallel, then reference-design(buck stage). No interface-spec agent.
\end{quote}
\paragraph*{P1}
\begin{quote}\small\ttfamily\raggedright
\# P1 Research - digest\\
~\\
- Roster: power-architect, scout(regulator), scout(powerpath), refdesign(buck).\\
  No interface-spec agent - no standards-bound interface on this board.\\
- Regulator: AP63356QZV-7 / AP63357QZV-7 (VDFN-13 2x3, sync, 3.5 A, 3.8-32 V,\\
  450 kHz, hiccup+OTP, \textasciitilde{}\$1.08@5). Family is ADJUSTABLE-ONLY (catalog "Fixed"\\
  is a scrape artifact). 5V/3A table: L 6.8 uH, Cin 10 uF, Cout 2x22 uF,\\
  Cbst 100 nF, FB 157k/30k. EN self-starts, SS internal 4 ms. refdesign\\
  prefers the PWM-only 356Q for the 50 mV ripple spec.\\
- power: 4 LAYERS REQUIRED, deciding corner 7 V low line. Eff 92.6/93.6/93.6 \%\\
  at 7/12/18 V, loss 1.03-1.20 W, Iin 2.31 A. Cin RMS ripple peaks 1.51 A at\\
  Vin=10 V. X7R only, no electrolytics. GND on In1+In2.\\
- powerpath: PMOS AOD403/AO4407A both \textless{}100 mW at 2.4 A. TVS floor SMBJ20A\\
  (32.4 V clamp). Inductors single-OEM (CENKER/XR). Terminals 14.1 mm vs 15 cap.\\
- CONFLICT for P2: thermal model used generic classes, never saw AP6335x\\
  (datasheet 25 C/W vs repo calibrated 51 C/W 4L) - run check\_thermal for real.\\
- Held for H1: fuse 4 vs 5 A, load transient/polymer cap, Tj\textless{}=105 C hard?
\end{quote}
\paragraph*{P2}
\emph{(no digest recorded)}
\subsection*{Gate attempt history}
\emph{(no entries)}
\subsection*{Issues}
\emph{(no entries)}
\subsection*{Human checkpoints}
{\small
\begin{longtable}{lllp{2.6cm}p{6.4cm}}
\toprule
\textbf{Id} & \textbf{Phase} & \textbf{Status} & \textbf{When} & \textbf{Note} \\
\midrule
\endhead
1 & P2 & not reached &  &  \\
2 & P4 & not reached &  &  \\
3 & P6 & not reached &  &  \\
4 & P8 & not reached &  &  \\
5 & P10 & not reached &  &  \\
\bottomrule
\end{longtable}
}
\section{Artifact Index}
{\small
\begin{longtable}{p{9cm}rp{4cm}}
\toprule
\textbf{File} & \textbf{Size} & \textbf{Note} \\
\midrule
\endhead
\texttt{state.json} & 7,532 B &  \\
\texttt{brief/brief.md} & 602 B &  \\
\texttt{requirements.md} & 14,844 B &  \\
\texttt{architecture/blocks.md} & 10,148 B &  \\
\texttt{architecture/constraints.json} & 7,534 B &  \\
\texttt{architecture/decisions.md} & 18,767 B &  \\
\texttt{architecture/power\_tree.md} & 10,096 B &  \\
\texttt{architecture/sheets.md} & 7,584 B &  \\
\texttt{architecture/stackup.md} & 14,439 B &  \\
\bottomrule
\end{longtable}
}
\end{document}
